\section{Metodologia Geral}
\label{sec:metodologia}

\subsection{Conjuntos de dados}

\paragraph{DOMPI-2025 (Questão 01).} O corpus DOMPI-2025~\cite{dataset_dompi2025}
reúne publicações do Diário Oficial dos Municípios do Piauí (DOM-PI) referentes a
2025, distribuídas em 12 territórios de desenvolvimento do estado. O texto foi
extraído dos PDFs originais com PaddleOCR, Docling e PyMuPDF, totalizando 77.337
documentos brutos. Trata-se de linguagem jurídico-administrativa: portarias,
decretos, extratos de contrato, editais de licitação e relatórios fiscais.

\paragraph{docentesDC (Questões 02 e 03).} O dataset
docentesDC~\cite{dataset_docentesdc} contém materiais acadêmicos (slides, notas de
aula, fragmentos de código e artigos) de 19 professores do Departamento de
Computação da UFPI, em 13.762 registros com os campos \texttt{text} e
\texttt{nome\_professor}. O conteúdo é heterogêneo e ruidoso: além de texto em
português extraído de slides, há artigos científicos em inglês fatiados em blocos e
código-fonte em linguagem C. É a base do pipeline de geração de pares descrito na
Seção~\ref{sec:q2}.

\subsection{Modelo e ambiente computacional}

Todos os experimentos partem do mesmo modelo base, \texttt{Qwen2.5-0.5B}, carregado
via biblioteca \textit{Transformers}~\cite{wolf2020transformers}. A execução ocorreu
em uma única sessão da plataforma Kaggle, em um \textit{notebook} unificado que
encadeia as três questões, com uma única GPU NVIDIA T4 (16~GB) visível. A escolha de
uma única T4 (em vez das duas disponíveis) foi deliberada: com duas GPUs, o
\texttt{Trainer} ativaria \textit{DataParallel} e concentraria os \textit{logits} do
vocabulário de 151~mil tokens em uma só placa, estourando a memória.

Por exigência do \textit{GradScaler} usado no treino em precisão mista
(\texttt{fp16} com \textit{autocast}), os pesos mestres dos treinos completos foram
mantidos em \texttt{fp32}; a T4 não possui suporte de hardware a \texttt{bfloat16}.
Empregou-se \textit{gradient checkpointing} em todos os treinos para reduzir o
consumo de memória, e o melhor \textit{checkpoint} de cada treino foi selecionado
pela menor perda de validação (\texttt{load\_best\_model\_at\_end}).

\subsection{Protocolo de avaliação}

Para cada experimento, o modelo é avaliado \emph{antes} e \emph{depois} do treino
sobre o mesmo conjunto reservado, usando as métricas da Seção~\ref{sec:metricas}. Na
Questão 01 as métricas cobrem a sequência completa (modelagem de linguagem pura);
nas Questões 02 e 03, apenas os tokens de resposta, conforme a
Seção~\ref{sec:mascara}. A avaliação quantitativa é complementada por comparações
qualitativas com geração determinística (\textit{greedy}), garantindo que as
respostas registradas antes e depois sejam reprodutíveis. Sementes aleatórias fixas
controlam embaralhamentos e divisões de dados. A Tabela~\ref{tab:visao-geral} resume
os quatro experimentos executados.

\begin{table}[ht]
  \centering
  \caption{Visão geral dos experimentos executados na sessão única de T4.}
  \label{tab:visao-geral}
  \footnotesize
  \begin{tabular}{@{}l l l r r r@{}}
    \toprule
    \textbf{Experimento} & \textbf{Técnica} & \textbf{Dataset} & \textbf{Params. treinados} & \textbf{Tempo} & \textbf{VRAM} \\
    \midrule
    Questão 01 & Pré-treino continuado & DOMPI-2025 & 494M (100\%) & 4h23 & 10,60~GB \\
    Questão 02 & SFT completo          & docentesDC & 494M (100\%) & 11~min & 9,77~GB \\
    Questão 03 & QLoRA (NF4)           & docentesDC & 2,16M (0,44\%) & 18~min & 1,77~GB \\
    Adicional  & QLoRA no 1.5B         & docentesDC & 4,36M (0,28\%) & 26~min & 2,75~GB \\
    \bottomrule
  \end{tabular}
\end{table}

\subsection{Ambientes e recursos das Questões 04 a 06}
\label{sec:amb-q456}

As Questões 04 a 06 ampliam o estudo para técnicas que não se resumem ao treino de
um único modelo e, por isso, foram desenvolvidas em ambientes próprios, com modelos e
ferramentas distintos dos das três primeiras. A Questão 04 usa dois modelos da
família Qwen2.5 na variante \texttt{-Instruct} (professor de 3 bilhões e aluno de 1,5
bilhão de parâmetros) e o mesmo corpus DOMPI-2025 da Questão 01, agora como fonte de
um dataset sintético de destilação. A Questão 05 não treina modelo algum: monta um
pipeline de RAG sobre o docentesDC com serviços da OpenAI para geração e
\textit{embeddings}. A Questão 06 reaproveita o próprio modelo pós-treinado na
Questão 02, adicionando-lhe camadas de \textit{guardrails} baseadas em regras. A
Tabela~\ref{tab:amb-q456} resume esses ambientes.

\begin{table}[ht]
  \centering
  \caption{Modelos, dados, ferramentas e ambientes das Questões 04 a 06.}
  \label{tab:amb-q456}
  \footnotesize
  \begin{tabularx}{\textwidth}{@{}l
    >{\hsize=1.05\hsize\raggedright\arraybackslash}X
    >{\hsize=0.85\hsize\raggedright\arraybackslash}X
    >{\hsize=1.30\hsize\raggedright\arraybackslash}X
    >{\hsize=0.80\hsize\raggedright\arraybackslash}X@{}}
    \toprule
    \textbf{Questão} & \textbf{Modelo(s)} & \textbf{Dataset} & \textbf{Principais ferramentas} & \textbf{Ambiente} \\
    \midrule
    Q4 Destilação & Qwen2.5-3B $\rightarrow$ 1.5B-Instruct & DOMPI-2025 (sintético, 100 trechos) & Transformers, PEFT/LoRA, SFTTrainer, LLM-as-Judge & GPU RTX 4050 (6~GB) \\
    Q5 RAG & \texttt{gpt-4o-mini} \newline + \texttt{text-\allowbreak embedding-\allowbreak 3-\allowbreak small} & docentesDC & LangChain, ChromaDB, Ragas & API OpenAI (local) \\
    Q6 Guardrails & Qwen2.5-0.5B (SFT da Q2) & docentesDC & Transformers + trilhos de regras (\textit{regex}) & CPU (local) \\
    \bottomrule
  \end{tabularx}
\end{table}
